% % % % % % % % % % % % % % % % % % % % % % % % % % %
% IS&T Template
% % % % % % % % % % % % % % % % % % % % % % % % % % %

%%%%%%%%%%%%%%%%%%%%%%%%%%%%%%%%%%
% Document class
%%%%%%%%%%%%%%%%%%%%%%%%%%%%%%%%%%
\documentclass[letterpaper,twocolumn,fleqn]{article}

%%%%%%%%%%%%%%%%%%%%%%%%%%%%%%%%%%
% Packages
%%%%%%%%%%%%%%%%%%%%%%%%%%%%%%%%%%
\usepackage{ist}
\usepackage[margin=1in]{geometry}
\usepackage{float}
\usepackage{wrapfig}
\usepackage{amsmath}
\usepackage{amssymb}
\usepackage{graphicx}
\usepackage{listings}
\usepackage{xcolor}
\usepackage{hyperref}
\usepackage{booktabs}
\pagestyle{empty}
\lstset{
  basicstyle=\ttfamily\footnotesize,
  breaklines=true,
  breakatwhitespace=true,
  columns=fullflexible,
  keepspaces=true,
  frame=single,
  showstringspaces=false,
  xleftmargin=0.02\columnwidth,
  xrightmargin=0.02\columnwidth
}

%%%%%%%%%%%%%%%%%%%%%%%%%%%%%%%%%%
% Title and Authors
%%%%%%%%%%%%%%%%%%%%%%%%%%%%%%%%%%

\title{Solving Linear Systems of Equations:\\
  Gaussian Elimination, Back Substitution, \& Cubic Splines}
\author{ Jared Acosta, Josiah Wiggins, Jake Davila, Jack Gaon, Jacky Li, Virgil Nhieu, Elijah Perez, Richie Wong\\ Cal Poly Pomona; Pomona, California \\ April 13, 2026}

%%%%%%%%%%%%%%%%%%%%%%%%%%%%%%%%%%
% Begin document
%%%%%%%%%%%%%%%%%%%%%%%%%%%%%%%%%%
\begin{document}
\maketitle

%%%%%%%%%%%%%%%%%%%%%%%%%%%%%%%%%%
\section{Introduction}
%%%%%%%%%%%%%%%%%%%%%%%%%%%%%%%%%%
Solving systems of linear equations $\mathbf{A}\mathbf{x} = \mathbf{b}$
is one of the most fundamental tasks in numerical computing.  Applications
range from circuit analysis and structural engineering to machine learning
and computer graphics.

This report presents Python implementations of \textbf{Gaussian
elimination with scaled partial pivoting} and \textbf{natural cubic
spline interpolation}, both translated from MATLAB code provided in
ECE~5110.  The original MATLAB \texttt{function\_solve\_lsoe} was
missing the back-substitution step; the \texttt{function\_splines} code
builds a constraint matrix and solves it via \texttt{solve\_lsoe}.  We
complete both algorithms in Python, verify them, and analyze accuracy.

%%%%%%%%%%%%%%%%%%%%%%%%%%%%%%%%%%
\section{Theory}
%%%%%%%%%%%%%%%%%%%%%%%%%%%%%%%%%%

\subsection{The Linear System}
Given an $n \times n$ matrix $\mathbf{A}$ and an $n \times 1$ vector
$\mathbf{b}$, we seek $\mathbf{x}$ such that
\[
\mathbf{A}\mathbf{x} = \mathbf{b}
\]

\subsection{Gaussian Elimination}
The algorithm transforms $\mathbf{A}$ into an upper-triangular matrix
$\mathbf{U}$ through a sequence of elementary row operations:

\textbf{Step 1 — Row Normalization.}
Scale each active row by its maximum absolute entry to prevent large
entries from dominating the pivot selection:
\[
\mathbf{A}_{r,:} \leftarrow \frac{\mathbf{A}_{r,:}}
                                  {\max_j |A_{r,j}|}, \quad
b_r \leftarrow \frac{b_r}{\max_j |A_{r,j}|}
\]

\textbf{Step 2 — Partial Pivoting.}
Swap rows so that the entry with the largest absolute value in the
current column becomes the pivot.  This reduces round-off error
amplification:
\[
\text{pivotRow} = \arg\max_{r \ge k} |A_{r,k}|
\]

\textbf{Step 3 — Forward Elimination.}
For each row $r > k$, subtract a multiple of the pivot row to zero
out the entry below the diagonal:
\[
m_{r,k} = \frac{A_{r,k}}{A_{k,k}}, \quad
\mathbf{A}_{r,:} \leftarrow \mathbf{A}_{r,:} - m_{r,k} \,\mathbf{A}_{k,:}
\]
After $n-1$ passes, $\mathbf{A}$ is upper-triangular.

\subsection{Back Substitution}
Once the system is in upper-triangular form $\mathbf{U}\mathbf{x} = \mathbf{c}$,
solve from the last row upward:
\[
x_i = \frac{1}{U_{i,i}}\!\left(c_i - \sum_{j=i+1}^{n} U_{i,j}\,x_j\right),
\quad i = n, n{-}1, \ldots, 1
\]
This is the step that was \textbf{missing} from the original MATLAB code.

\subsection{Natural Cubic Splines}
Given $n$ data points $(x_0, y_0), \ldots, (x_{n-1}, y_{n-1})$,
we fit $n-1$ cubic polynomials
$S_i(x) = a_i x^3 + b_i x^2 + c_i x + d_i$
with \textbf{five types of constraints}:
\begin{enumerate}
  \item Each $S_i$ passes through its left knot: $S_i(x_i) = y_i$
  \item Each $S_i$ passes through its right knot: $S_i(x_{i+1}) = y_{i+1}$
  \item First-derivative continuity: $S_i'(x_{i+1}) = S_{i+1}'(x_{i+1})$
  \item Second-derivative continuity: $S_i''(x_{i+1}) = S_{i+1}''(x_{i+1})$
  \item Natural BCs: $S_0''(x_0) = 0$, $S_{n-2}''(x_{n-1}) = 0$
\end{enumerate}
This yields $4(n-1)$ equations for $4(n-1)$ unknowns, which we solve
using \texttt{solve\_lsoe}.

%%%%%%%%%%%%%%%%%%%%%%%%%%%%%%%%%%
\section{Implementation}
%%%%%%%%%%%%%%%%%%%%%%%%%%%%%%%%%%

\subsection{solve\_lsoe (tools.py)}
The complete method in \texttt{lib/tools.py}:

\begin{lstlisting}
def solve_lsoe(self, A, B):
    A = np.array(A, dtype=np.float64)
    B = np.array(B, dtype=np.float64).flatten()
    n = A.shape[0]

    # Forward Elimination
    for ix in range(n - 1):
      # 1. NORMALIZE
      for r in range(ix, n):
        max_row = np.max(np.abs(A[r, ix:]))
        if max_row < 1e-15: continue
        A[r,ix:] = A[r,ix:] / max_row
        B[r] = B[r] / max_row

      # 2. PARTIAL PIVOTING
      pivot = ix + np.argmax(np.abs(A[ix:,ix]))
      if pivot != ix:
        A[[ix,pivot]] = A[[pivot,ix]]
        B[[ix,pivot]] = B[[pivot,ix]]

      # 3. ELIMINATE
      for r in range(ix+1, n):
        factor = A[r,ix] / A[ix,ix]
        A[r,ix:] -= factor * A[ix,ix:]
        B[r] -= factor * B[ix]

    # 4. BACK SUBSTITUTION
    x = np.zeros(n)
    for i in range(n-1, -1, -1):
      x[i] = (B[i] - A[i,i+1:] @ x[i+1:])
              / A[i,i]
    return x, 0
\end{lstlisting}

%%%%%%%%%%%%%%%%%%%%%%%%%%%%%%%%%%
\section{Test Cases}
%%%%%%%%%%%%%%%%%%%%%%%%%%%%%%%%%%

\subsection{Instructor's 4×4 System}
From the MATLAB \texttt{unittest\_lsoe.m}:
\[
\mathbf{A} = \begin{bmatrix}
1 & 1 & 1 & 1 \\
1 & 2 & 3 & 4 \\
0 & 2 & 0 & 1 \\
-1 & -2 & 1 & 2
\end{bmatrix}, \quad
\mathbf{b} = \begin{bmatrix} 4 \\ 10 \\ 12 \\ 0 \end{bmatrix}
\]
Our solver produces $\mathbf{x} = [5.5,\; 1,\; {-12.5},\; 10]^T$,
matching \texttt{numpy.linalg.solve} to within $1.78 \times 10^{-15}$.

\subsection{Additional Tests}
\begin{center}
\begin{tabular}{lcc}
\toprule
\textbf{Test} & \textbf{Size} & \textbf{Result} \\
\midrule
Identity matrix        & $3 \times 3$ & $\checkmark$ \\
Simple 2$\times$2      & $2 \times 2$ & $\checkmark$ \\
Singular matrix        & $3 \times 3$ & detected \\
Random well-conditioned & $6 \times 6$ & $\checkmark$ \\
Cubic spline (sin, 6 knots) & $20 \times 20$ & $\checkmark$ \\
Cubic spline (4-point)   & $12 \times 12$ & $\checkmark$ \\
\bottomrule
\end{tabular}
\end{center}

%%%%%%%%%%%%%%%%%%%%%%%%%%%%%%%%%%
\section{Results}
%%%%%%%%%%%%%%%%%%%%%%%%%%%%%%%%%%
Figure~\ref{fig:acc} shows the accuracy of our solver as a function
of matrix size $n$ (left panel) and the condition number of the
randomly generated test matrices (right panel).

\begin{figure}[H]
\centering
\includegraphics[width=\columnwidth]{unit04_accuracy_analysis.png}
\caption{Left: maximum absolute error vs.\ matrix size.
Right: condition number $\kappa(\mathbf{A})$ vs.\ matrix size.
All test matrices are diagonally dominant with $\mathbf{A} = \text{randn}(n) + n\,\mathbf{I}$.}
\label{fig:acc}
\end{figure}

The error remains near machine precision ($\sim 10^{-13}$) for
well-conditioned systems, confirming that our implementation is
numerically stable for diagonally dominant matrices.

Figure~\ref{fig:spline} shows the cubic spline interpolation of
$\sin(x)$ with 6 knots.

\begin{figure}[H]
\centering
\includegraphics[width=\columnwidth]{unit04_cubic_spline.png}
\caption{Left: cubic spline fit of $\sin(x)$ using 6 knots.
Right: absolute interpolation error.}
\label{fig:spline}
\end{figure}

%%%%%%%%%%%%%%%%%%%%%%%%%%%%%%%%%%
\section{Discussion}
%%%%%%%%%%%%%%%%%%%%%%%%%%%%%%%%%%
\begin{itemize}
  \item \textbf{Scaled partial pivoting} prevents catastrophic
        cancellation when rows have vastly different magnitudes.
  \item \textbf{Back substitution} is $O(n^2)$ and dominates only
        for small $n$; the overall complexity is $O(n^3)$ due to
        forward elimination.
  \item For ill-conditioned systems ($\kappa \gg 1$), even partial
        pivoting cannot prevent loss of significant digits.
        Full pivoting or iterative refinement would be needed.
  \item Our implementation detects singular matrices by checking
        whether any pivot falls below $10^{-15}$.
\end{itemize}

%%%%%%%%%%%%%%%%%%%%%%%%%%%%%%%%%%
\section{Conclusion}
%%%%%%%%%%%%%%%%%%%%%%%%%%%%%%%%%%
We translated two MATLAB functions to Python: \texttt{solve\_lsoe}
(with the missing back substitution added) and \texttt{function\_splines}
(natural cubic spline interpolation).  The spline method builds a
$4(n{-}1) \times 4(n{-}1)$ constraint system and solves it directly
via our Gaussian elimination solver.  All seven test cases pass
with errors at machine-precision level.

\end{document}
